\documentclass[11pt]{article}
\usepackage{amsmath,amsfonts,amssymb,amsthm}
\usepackage{graphicx}
\usepackage{algorithm}
\usepackage{algorithmic}
\usepackage{float}
\usepackage{hyperref}
\usepackage{xcolor}
\usepackage{booktabs}

\title{Integral Representations of Sobolev Spaces via ReLU$^k$ Activation Function and Optimal Error Estimates for Linearized Networks}
\author{Xinliang Liu, Tong Mao, Jinchao Xu}
\date{}

\begin{document}
\maketitle

\section{Integral Representations of Sobolev Spaces via ReLU$^k$ Activation Function and Optimal Error Estimates for Linearized Networks}

\subsection{Introduction and Background}

This paper presents fundamental theoretical results about shallow neural networks with ReLU$^k$ activation functions. The work focuses on linearized shallow networks, where inner parameters are fixed and only linear coefficients are optimized, demonstrating they can achieve optimal approximation rates in Sobolev spaces.

The standard shallow ReLU$^k$ neural network function class is defined as:
\begin{equation}
\Sigma_n^k = \left\{\sum_{j=1}^n a_j \sigma_k(\theta_j \cdot \tilde{x}) : \theta_j \in \mathbb{S}^d, a_j \in \mathbb{R} \right\}
\end{equation}
where $\tilde{x} = \begin{pmatrix} x \\ 1 \end{pmatrix}$ and $\theta_j = \begin{pmatrix} w_j \\ b_j \end{pmatrix}$. The stable version with bounded parameters is:
\begin{equation}
\Sigma_{n,M}^k = \left\{\sum_{j=1}^n a_j \sigma_k(\theta_j \cdot \tilde{x}) : \theta_j \in \mathbb{S}^d, \sum_{j=1}^n |a_j| \leq M \right\}
\end{equation}

The paper introduces and analyzes the finite neuron space (FNS), a linear subset of $\Sigma_n^k$:
\begin{equation}
L_n^k = L_n^k(\{\theta_j^*\}_{j=1}^n) = \text{span} \{\phi_1, \ldots, \phi_n\}
\end{equation}
where $\phi_j(x) = \sigma_k(\theta_j^* \cdot \tilde{x})$ with predetermined parameters $\{\theta_j^*\}_{j=1}^n \subset \mathbb{S}^d$. The constrained version is:
\begin{equation}
L_{n,M}^k = \left\{\sum_{j=1}^n a_j \phi_j : \left(\sum_{j=1}^n a_j^2\right)^{1/2} \leq M \right\}
\end{equation}

\subsection{Main Theoretical Results}

The paper establishes three key theoretical results:

\subsubsection{Novel Integral Representation of Sobolev Spaces}

The paper proves that Sobolev spaces $H^{d+2k+1/2}(\Omega)$ can be characterized through an integral representation using ReLU$^k$ functions:
\begin{equation}
H^{d+2k+1/2}(\Omega) = \left\{\int_{\mathbb{S}^d} \sigma_k(\theta \cdot \tilde{x}) \psi(\theta) d\theta : \psi \in L^2(\mathbb{S}^d) \right\}
\end{equation}

This representation provides a direct link between Sobolev spaces and neural networks, showing that for any $f \in H^{d+2k+1/2}(\Omega)$, there exists $\psi \in L^2(\mathbb{S}^d)$ such that:
\begin{equation}
f(x) = \int_{\mathbb{S}^d} \sigma_k(\theta \cdot \tilde{x}) \psi(\theta) d\theta
\end{equation}

with the norm equivalence:
\begin{equation}
\|f\|_{H^{d+2k+1/2}(\Omega)} \simeq \inf_{\psi \in L^2(\mathbb{S}^d)} \{\|\psi\|_{L^2(\mathbb{S}^d)} : f(x) = \int_{\mathbb{S}^d} \sigma_k(\theta \cdot \tilde{x}) \psi(\theta) d\theta\}
\end{equation}

\subsubsection{Optimal Approximation Rates for Linearized Networks}

The paper demonstrates that linearized shallow networks with well-distributed parameters achieve optimal approximation rates. Specifically, for well-distributed parameters $\{\theta_j^*\}_{j=1}^n \subset \mathbb{S}^d$ and any $f \in H^{d+2k+1/2}(\Omega)$, with $M \simeq \|f\|_{H^{d+2k+1/2}(\Omega)}$:
\begin{equation}
\inf_{f_n \in L_{n,M}^k} \|f - f_n\|_{L^2(\Omega)} \lesssim n^{-1/2-2k+1/2d} \|f\|_{H^{d+2k+1/2}(\Omega)}
\end{equation}

A crucial finding is that this rate matches the optimal approximation rate for nonlinear shallow networks (up to logarithmic factors), showing linearized networks are as powerful as their nonlinear counterparts for functions in appropriate Sobolev spaces.

\subsubsection{Coefficient Bound Estimates}

The paper establishes the important coefficient estimate:
\begin{equation}
\sqrt{n}\|a\|_2 \lesssim \|f\|_{H^{d+2k+1/2}(\Omega)}
\end{equation}

This bound is essential for practical applications, particularly in generalization analysis, as it ensures numerical stability and provides control over the complexity of the approximating functions.

\subsection{Algorithmic Approach}

The paper focuses on a linearized approach to neural network approximation, which can be summarized algorithmically:

\begin{enumerate}
    \item Select a set of well-distributed parameters $\{\theta_j^*\}_{j=1}^n \subset \mathbb{S}^d$, where the mesh size $h = \max_{\theta \in \mathbb{S}^d} \min_{1 \leq j \leq n} \rho(\theta, \theta_j^*)$ satisfies $h \lesssim n^{-1/d}$
    
    \item For these fixed parameters, construct the finite neuron space:
    \begin{equation}
    L_n^k = \text{span}\{\sigma_k(\theta_1^* \cdot \tilde{x}), \ldots, \sigma_k(\theta_n^* \cdot \tilde{x})\}
    \end{equation}
    
    \item For a target function $f \in H^{d+2k+1/2}(\Omega)$, find the approximation $f_n \in L_{n,M}^k$ by solving the convex least squares problem:
    \begin{equation}
    f_n = \arg\min_{g \in L_{n,M}^k} \|f - g\|_{L^2(\Omega)}
    \end{equation}
    where $M \simeq \|f\|_{H^{d+2k+1/2}(\Omega)}$
\end{enumerate}

This approach transforms the challenging non-convex optimization problem of training full neural networks into a simpler convex optimization problem that can be solved efficiently and reliably.

\subsection{Numerical Results and Practical Implications}

The paper applies the theoretical results to analyze several practical problems:

\subsubsection{Application to Random Feature Methods}

The paper shows that randomized approaches (where parameters $\{\theta_j\}$ are randomly sampled) can be understood through the lens of the deterministic theory. If $\{\theta_j\}_{j=1}^n$ are i.i.d. uniform samples from $\mathbb{S}^d$, the expected error satisfies:
\begin{equation}
\mathbb{E}_n\left[\inf_{a \in \mathbb{R}^n} \left\|f - \sum_{j=1}^n a_j \sigma_k(\theta_j \cdot \tilde{x})\right\|_{L^2(\Omega)}\right] \lesssim \left(\frac{n}{\log n}\right)^{-1/2-2k+1/2d} \|f\|_{H^{d+2k+1/2}(\Omega)}
\end{equation}

This is only slightly worse (by a logarithmic factor) than the deterministic approach with well-distributed points. The analysis suggests that randomization works primarily because random points tend to be well-distributed with high probability, not because of any inherent advantage of randomness.

\subsubsection{Generalization Analysis for PDEs}

The paper presents generalization error bounds for learning solutions to elliptic PDEs using the FNS. For a second-order linear elliptic PDE:
\begin{align}
-\Delta g(x) + g(x) &= h(x), \quad x \in \Omega \\
\frac{\partial g}{\partial n}(x) &= 0, \quad x \in \partial\Omega
\end{align}

When approximating the solution using $m$ training samples, the expected $H^1$ error satisfies:
\begin{equation}
\mathbb{E}_{x_1,\ldots,x_m}[\|f_{n,m} - f\|^2_{H^1(\Omega)}] \lesssim (\|h\|_{L^\infty(\Omega)} + M)M m^{-1/2}
\end{equation}

where $f_{n,m}$ is the empirical risk minimizer in $L_{n,M}^k$ and $f$ is the true solution.

\subsubsection{Comparison with Traditional Approximation Methods}

The paper proves that for functions in $H^{d+2k+1/2}(\Omega)$, the FNS can outperform classical approximation methods like finite elements. Specifically, for the optimal finite element space $V_n^k$ (piecewise polynomials of degree $\leq k$ on $n$ elements):
\begin{equation}
\inf_{f_n \in L_n^k} \|f - f_n\|_{L^2(\Omega)} = O\left(n^{-1/2(1-1/d)}\right) \inf_{f_n \in V_n^k} \|f - f_n\|_{L^2(\Omega)}
\end{equation}

This demonstrates that linearized neural networks can overcome the curse of dimensionality that affects traditional finite element methods.

\subsection{Mathematical Techniques and Analysis}

The paper employs several sophisticated mathematical techniques:

\begin{enumerate}
    \item Spherical harmonic analysis and Legendre polynomial expansions to analyze functions on the sphere
    
    \item Careful analysis of the Legendre coefficients of ReLU$^k$ functions
    
    \item Construction of specialized operators that map between functions on the sphere and functions on the ball while preserving the structure of ReLU$^k$ functions
    
    \item Analysis of matrices induced by Legendre polynomials evaluated at scattered points on the sphere
    
    \item Use of Rademacher complexity to bound generalization error
\end{enumerate}

These techniques allow the authors to establish sharp approximation rates and precise coefficient bounds, providing a rigorous mathematical foundation for linearized neural networks.

\subsection{Practical Significance}

The findings have several significant practical implications:

\begin{enumerate}
    \item Linear training of neural networks with fixed inner weights can achieve optimal approximation rates, potentially offering dramatic computational advantages over full network training
    
    \item The optimal approximation rates are achieved by well-distributed (deterministic) points, suggesting deterministic parameter selection could be preferable to random initialization
    
    \item The coefficient bounds provide guarantees for numerical stability and generalization capability
    
    \item The approach can be particularly effective for solving PDEs, offering advantages over traditional numerical methods in high dimensions
\end{enumerate}

These results help explain why methods like extreme learning machines, random feature methods, and randomized neural networks have been successful in practice, while suggesting potential improvements through deterministic parameter selection.

\end{document} 